\section{Related Work}
\label{sec:related}

\subsection{Information-theoretic bounds on traffic analysis}
\label{sec:related/bounds}

The idea of bounding what traffic analysis can achieve, rather than measuring
what a particular analyzer does achieve, is well established, and the results
split by the structure they assume. \citet{cherubin2017bayes} gives
feature-set-conditional Bayes-error bounds for website fingerprinting, which is
a closed-set floor and the closest precedent for the upper side of our bracket;
its bound is conditional on the feature set the analyst chooses, so it answers
how well one can do with given features rather than how much separation the
observable carries. The anonymity and quantitative-information-flow literature
takes the complementary route, from an entropy-like quantity or a metric on
distributions toward a bound on inference
\citep{serjantov2003metric,ghaderi2010anonymous,venkitasubramaniam2008eavesdroppers,venkitasubramaniam2008chaum,chatzikokolakis2007error,chatzikokolakis2008noisy}.
Those bounds are typically evaluated on models or synthetic channels, which is
where our contribution differs: we measure an interval on captured traffic.

The nearest recent results are the leakage-inevitability theorem of
\citet{liu2026inevitability} and the rate-distortion frontier of
\citet{liu2026rate}. Both are single-letter in a precise sense that matters for
this paper. Exponential error decay in the first requires conditionally i.i.d.\
repeated observations, and the frontier in the second is proved inside a
stationary-memoryless defense class. Packet prefixes are neither conditionally
independent nor stationary, so neither result applies to the object we measure
without an assumption that the data does not satisfy. Our upper side is
distribution-free precisely so that the argument does not depend on that
assumption; the price is that the lower side is asymptotic rather than
exponential.

\subsection{Total variation estimation and lower bounds}
\label{sec:related/tv}

Turning a radius into a usable statement requires finite-sample control of the
empirical distribution, and the Kolmogorov-Smirnov family supplies it
\citep{wei2011dkwtwo,reeve2024dkwmassart} with high-probability and two-sample
refinements that tighten the constant
\citep{diakonikolas2020highprob,han2015l1distribution,jiao2018l1distance,chan2013closeness}. Two
routes give lower bounds. Classifier-based methods convert a decision function's
error rate into a separation, and \citet{michel2020high} provides the adaptive
high-probability form we use. \citet{tao2024discriminative} connects
discriminative total-variation estimation to classification risk under a
hypothesis-class assumption, which is the route the conditional large-prefix
bound follows. Our construction composes these pieces into an explicit two-sided
interval and reports the two sides separately, because the composition is where
the guarantee asymmetry becomes visible.

\subsection{Open-set recognition and its comparators for encrypted traffic}
\label{sec:related/openset}

Open-set recognition is the framing this paper's setting inherits.
\citet{scheirer2013toward} introduced the risk of novel classes at deployment,
and generalization and learnability analyses
\citep{fang2021learning,garov2025closer} characterize what can be guaranteed for
detectors that must reject unknowns. Those results bound the risk of a learning
procedure; a bracket bounds the separation present in the data. The two are
complementary, since a learnability bound says how well one can train and a
bracket says how much room there is to train into. Work on unknown-traffic
identification in the encrypted setting
\citep{hu2022opencbd,li2023facing} applies the same rejection framing to network
data and reports detector-level metrics.

Two comparators bear directly on our positioning, and we state the evidential
basis for each. \citet{shamsimukhametov2023indistinguishability} studies
indistinguishability under open TLS parameters with Encrypted ClientHello; it is
the closest protocol-privacy precedent and the work our channel declaration must
answer to, since it addresses how much of the observable survives parameter
changes while we report an interval on the observable as declared, with the
identifier's contribution measured rather than assumed away. Our knowledge of
this work is abstract-level: its publisher record and abstract were retrieved
but the full text was not, and the novelty boundary in
\Cref{sec:related/novelty} is stated accordingly. RoNeTC
\citep{wang2025reliable} represents uncertainty-aware reliable open-set
classification, and a detector of that family is one arm of our validation
panel. Against these and against prototype-radius rejection
\citep{qu2026sepspace}, the distinction we draw is not that their detectors are
weak but that their reported quantities are accuracy-style metrics with no
stated limit, whereas a bracket says what limit applies.

\subsection{Data hygiene and shortcut learning}
\label{sec:related/hygiene}

Three works together motivate the evaluation discipline this paper enforces, and
we cite them as a trio because each covers a different failure mode.
\citet{zhao2025sweet} shows that representation-learning performance in
encrypted-traffic classification can be driven by shortcut features and leakage
in the data organization. \citet{wang2026bias} provides a model-agnostic,
data-driven framework for detecting dataset-specific shortcuts and
environment-entangled correlations. \citet{wickramasinghe2025sok} systematizes
this critique across the classifier literature, which is what makes it
load-bearing rather than redundant: it establishes that the problem is a
property of the field's evaluation practice rather than of one paper's setup.
A fourth line of work turns the critique into an audit procedure rather than a
diagnosis, randomizing address fields and auditing for shortcut-induced
unreliability at evaluation time \citep{zhao2026shortcutcatcher}; the
randomization we apply is the same instrument.
Together these are why we use flow- and time-disjoint splits, prohibit
packet-level splitting, randomize address fields, and verify the masking regime
instead of assuming it. The hygiene regime in \Cref{sec:design/channel} is a
direct response to these three, and the DST\_ASN audit we report is the same
instrument applied to our own environment confound.

\subsection{Unseen-class mass and protocol context}
\label{sec:related/mass}

A bracket over observed services says nothing about services absent from the
corpus, and the unseen-species literature supplies the machinery for that gap
\citep{orlitsky2016optimal,eriksson2026unseen}, which we use for the
deployment-site term in \Cref{sec:discussion/missing_mass}. On the protocol
side, QUIC's wire image and its operational implications are documented by
\citet{kuhlewind2022rfc}, the corpus we measure on and its companion are
described by \citet{cesnet2025cesnet} and \citet{luxemburk2023cesnet}, and prior
work establishes the difficulty of analysis in this setting
\citep{smith2021website,li2023facing,jerabek2025simple,luxemburk2023encryptedquic,zhang2023applicationquic}.
Learning-based traffic analyzers are now numerous
\citep{lin2022etbert,zhao2023yatc,meng2023netgpt,cui2025trafficllm,hu2022opencbd,athanasiou2025selfdefense,dyer2012peekaboo};
our contribution is orthogonal to their architectures, since it concerns what
any of them can achieve on a declared observable rather than how one of them is
built.

\paragraph{Novelty boundary.}
\label{sec:related/novelty}
To our knowledge, no prior work measures a two-sided bracket with asymmetric
guarantee strength on real encrypted traffic. We state this as a knowledge claim
rather than a search result, and the search behind it is recorded with its own
limits: it was a motivated search whose query strings and databases were not
captured at search time, so it is not a reproducible query log, and a motivated
search can miss work phrased in different vocabulary (for example
distinguishability certificates, or tolerance intervals on distributions over
observable QoS quantities). The comparison is against the works above and the
others surveyed, each of which reports a bound on synthetic or model channels, a
closed-set error floor, or an accuracy-style detector metric, and none of which
reports a per-pair interval on captured traffic whose sides carry explicitly
different certification strength. Our closest comparator,
\citeauthor{shamsimukhametov2023indistinguishability}, is known to us at
abstract level only, so the boundary it defines is stated at that strength
rather than strengthened by assumption.
